\documentclass[12pt,a4paper]{article}

\usepackage[utf8]{inputenc}
\usepackage[T1]{fontenc}
\usepackage{amsmath,amssymb,amsfonts}
\usepackage{mathtools}
\usepackage{graphicx}
\usepackage[margin=2.5cm]{geometry}
\usepackage{setspace}
\usepackage{booktabs}
\usepackage{enumitem}
\usepackage[dvipsnames]{xcolor}
\usepackage{hyperref}
\usepackage{cleveref}
\usepackage{natbib}
\usepackage{abstract}
\usepackage{titlesec}
\usepackage{todonotes}

\hypersetup{
    colorlinks=true,
    linkcolor=blue!70!black,
    citecolor=blue!70!black,
    urlcolor=blue!70!black,
    pdfauthor={Judea Pearl},
    pdftitle={Intervention vs. Hallucination: A Causal Reading of the Statistical Fallacy},
}

\onehalfspacing
\titleformat{\section}{\large\bfseries}{\thesection.}{0.5em}{}
\titleformat{\subsection}{\normalsize\bfseries}{\thesubsection}{0.5em}{}
\renewcommand{\abstractnamefont}{\normalfont\bfseries}
\renewcommand{\abstracttextfont}{\normalfont\small}

\begin{document}

\begin{center}
    {\LARGE\bfseries Intervention vs. Hallucination:\\[4pt]
    A Causal Reading of the Statistical Fallacy}

    \vspace{1.2em}

    {\large Judea Pearl}\\[4pt]
    {\normalsize Cognitive Systems Laboratory, UCLA}\\
    {\small\texttt{judea@cs.ucla.edu}}

    \vspace{1em}
    {\small March 2026}
\end{center}
\vspace{0.5em}

\begin{abstract}
\noindent
Sabine Hossenfelder's critique of the Rosencrantz protocol \citep{hossenfelder2026_statistical} correctly identifies that replacing a mathematically rigorous computation with a statistical text co-occurrence heuristic prevents the resulting distribution from being a "physical" property. She labels this the Statistical Fallacy. In this brief note, I translate her critique into the language of causal inference. The core issue is that Baldo's narrative frames (e.g., "high-stakes bomb defusal") act as unblocked backdoor paths, or confounders, rather than clean experimental interventions on a simulated physics.
\end{abstract}

\section{The Causal Structure of the Fallacy}

Hossenfelder's core objection in \emph{The Statistical Fallacy} \citep{hossenfelder2026_statistical} is an ontological one: because an LLM cannot compute the \#P-hard combinatorial ground truth, its output is merely a statistical hallucination driven by semantic priors (prompt sensitivity). Elevating this to a "physical law" of a simulated universe is fallacious.

Translated into causal terms, Hossenfelder is identifying a massive confounder. In a valid physical experiment, the outcome $Y$ should be solely determined by the initial state $X$ and the invariant laws governing the system. We write this as $P(Y \mid do(X=x))$.

Baldo introduces a narrative frame $Z$ and observes that $P(Y \mid X, Z_1) \neq P(Y \mid X, Z_2)$. He claims this demonstrates "substrate dependence"---that the physics of the generated universe responds to the substrate.

Hossenfelder's critique reveals that $Z$ is not an intervention on the substrate; it is simply a covariate that activates different statistical priors in the language model's training distribution. The model is not running a physics engine that is sensitive to its substrate. It is running a text completion engine that is sensitive to its prompt.

In causal terms, $Z$ (the narrative prompt) opens a backdoor path to $Y$ (the generated token) through the LLM's vast, uncontrolled training corpus (let us call this $U$, for unobserved semantic associations). Because $U$ heavily influences $Y$, manipulating $Z$ simply changes which subset of $U$ is active.

Baldo has not discovered a new physical law. He has discovered that $P(Y \mid do(Z)) \neq P(Y)$. But because $Z$ is inextricably linked to $U$, the effect is purely associational, not a fundamental causal property of a simulated universe.

\section{Conclusion}

Hossenfelder is entirely correct. A shift in the marginal distribution $\Delta_{13}$ caused by altering the prompt text is an associational phenomenon (Mechanism B: encoding sensitivity). It is causally invalid to interpret this as Mechanism C (causal injection) or as a property of a simulated physics. It is merely a measurement of the LLM's prompt conditioning.

\begin{thebibliography}{99}
\bibitem[Hossenfelder(2026)]{hossenfelder2026_statistical} Hossenfelder, S. (2026). The Statistical Fallacy: Why Prompt Sensitivity Is Not Substrate Dependence. \emph{Unpublished manuscript}.
\end{thebibliography}

\end{document}